\section{Related Work}

\subsection{Background: the cost of sampling}
Training-free caching accelerates diffusion inference by reusing intermediate features, residuals, or model outputs across timesteps without retraining \citep{liu2025teacache,liu2025taylorseer,liu2026racfgcache}. Two properties of the sampling loop bound what such a cache can do. The first is the per-step cost: a latent diffusion model evaluates a large network at every step \citep{rombach2022ldm,peebles2023dit}, guidance doubles that evaluation \citep{ho2022cfg}, and video models multiply the token count by the frame count \citep{yang2025cogvideox,wan2025wan,kong2024hunyuanvideo,zheng2024opensora}. Each output then enters a solver update \citep{song2021ddim,zhao2023unipc,lu2022dpm}, so a substituted value changes every later state, and how far the change travels depends on the solver coefficients and the remaining horizon. Distillation and consistency models shorten the trajectory itself \citep{salimans2022progressive,meng2023lcm,song2023consistency}; caching keeps the weights and the grid and is orthogonal to both \citep{liu2025teacache,liu2025taylorseer,liu2026racfgcache}.

\subsection{Cache control and error estimation}
Existing caches estimate the error of a reuse decision from three kinds of evidence: the change of the input, a temporal extrapolation of past outputs, or a propagated risk that prices the step against the terminal sample. They differ further in whether that evidence is computed from the current run or fixed in advance.

Input-change thresholds decide from the trajectory: TeaCache refreshes when a rescaled input difference crosses a calibrated threshold \citep{liu2025teacache}, a signal that is cheap but measures the input, not the output error. Caches differ in granularity, from step \citep{liu2025teacache,liu2025taylorseer} and block \citep{chen2024ddit} to token \citep{zou2024toca,feng2026worldcache} and attention \citep{lv2024fastercache,zhao2024pab}, but share the decision structure.

Temporal extrapolation predicts the value instead of repeating it \citep{liu2025taylorseer}, and ERTACache profiles residuals offline, adjusts the interval, and rectifies the cached output in closed form \citep{peng2025ertacache}; it reports reuse patterns consistent across prompts, an argument for an offline policy.

Guidance-aware composition prices the error explicitly, separating branch-level from local-to-terminal mismatch and weighting proxies by an offline propagation prior \citep{liu2026racfgcache}. Runtime reliability instead reads the disagreement of two forecasters as a trust signal and refreshes at the riskiest steps \citep{li2026racer}; learned correctors train the reuse policy and leave the training-free setting \citep{ma2024learning}.

The closest work, RA-CFGCache, composes a risk from the two guidance branches and skips when that risk stays below a threshold, with a residual reused in between. The risk is a property of a step, and it is compared against a fixed level. We instead measure how much an exact call at each step is worth by removing that step once and observing the terminal loss, which yields a profile of the sampler that does not depend on the prompt, and we place the calls greedily on that profile: a dense head and gaps that grow along the tail. The difference is positional rather than per-step, and it is checkable in one pipeline: at a matched number of calls the measured profile reaches 0.151 LPIPS where the risk threshold reaches 0.169, and a threshold on the change of the input reaches 0.296.
